% Preamble untuk Proposal dan Skripsi - UIN Imam Bonjol Padang
\usepackage[indonesian]{babel}
\usepackage[utf8]{inputenc}
\usepackage[T1]{fontenc}
\usepackage{geometry}
\usepackage{setspace}
\usepackage{indentfirst}
\usepackage{graphicx}
\usepackage{mathptmx} % Times New Roman
\usepackage{titlesec} % Untuk format judul bab
\usepackage{caption} % Untuk format caption
\usepackage{etoolbox} % Untuk reset penomoran section per chapter
\usepackage{amsmath} % Untuk \numberwithin
\usepackage{lipsum} % Untuk dummy text (testing)
\usepackage{fancyhdr} % Untuk pengaturan posisi nomor halaman
\usepackage[hidelinks]{hyperref} % Untuk metadata PDF dan link
\usepackage{array} % Untuk format kolom tabel (p{...} dengan centering)
\usepackage[titles]{tocloft} % Untuk format Daftar Isi (dengan opsi titles agar tidak konflik dengan titlesec)
\usepackage{microtype} % Membantu perataan teks agar lebih rapi tanpa banyak hyphenation
\usepackage{enumitem} % Untuk pengaturan list (enumerate/itemize)
\usepackage{pgfgantt} % Untuk membuat Gantt Chart
\usepackage{float} % Untuk anchoring tabel/gambar dengan [H]
\usepackage{tikz} % Untuk membuat Flowchart
\usepackage{longtable} % Untuk tabel yang lebih dari satu halaman
\usetikzlibrary{shapes.geometric, arrows, positioning}

% --- Pengaturan Sitasi (APACite dengan Natbib Compatibility) ---
\usepackage[natbibapa]{apacite}

% Lokalisasi Sitasi (et al -> dkk, and -> dan)
\AtBeginDocument{
  \renewcommand{\BOthers}[1]{, dkk.\hbox{}}
  \renewcommand{\BOthersPeriod}[1]{, dkk.\hbox{}}
  \renewcommand{\BAnd}{dan}
  \renewcommand{\BRetrievedFrom}{Diambil dari}
  \renewcommand{\BRetrieved}[1]{Diakses pada #1}
}

% Menyesuaikan format agar selalu menggunakan kurung (Author, Year)
\let\oldcite\cite
\renewcommand{\cite}{\citep}

% --- Pengaturan List (Enumerate/Itemize) ---
% Menghapus spasi tambahan antar item dan sekitar list agar sesuai dengan spasi 1.5 paragraf
\setlist{nosep, leftmargin=1.25cm, labelsep=0.5cm}
\setlist[enumerate]{labelindent=\parindent}

% --- Pengaturan Hyphenation (Pemenggalan Kata) ---
% Mematikan pemenggalan kata otomatis
\hyphenpenalty=10000
\exhyphenpenalty=10000
\sloppy % Memberikan toleransi spasi antar kata agar tetap justify meskipun tanpa pemenggalan

% --- Variabel Global Identitas ---
\newcommand{\mhsname}{MUHAMAD IFFAN RAMADHAN}
\newcommand{\mhsnim}{2117020036}
\newcommand{\mhsjudul}{Analisis Performa dan Akurasi Model Whisper pada Sistem Transkripsi Terdistribusi Menggunakan Next.js dan Python Worker}
\newcommand{\mhsprodi}{Sistem Informasi}
\newcommand{\mhsfakultas}{Fakultas Sains dan Teknologi}
\newcommand{\mhsuniv}{Universitas Islam Negeri Imam Bonjol Padang}
\newcommand{\mhstahun}{2026}

% --- Variabel Global Pembimbing ---
\newcommand{\pembimbingI}{Ahmad Fauzi, M.T.I}
\newcommand{\nipI}{123456789}
\newcommand{\pembimbingII}{Ozzy Secio Riza, M.Kom}
\newcommand{\nipII}{987654321}

% --- Variabel Global Pejabat ---
\newcommand{\kaprodi}{Nama Ketua Program Studi}
\newcommand{\nipKaprodi}{19800101XXXXXXXX}
\newcommand{\dekan}{Nama Dekan Fakultas}
\newcommand{\nipDekan}{19700101XXXXXXXX}

% Pengaturan Margin
\geometry{
  a4paper,
  top=4cm,
  bottom=3cm,
  left=4cm,
  right=3cm
}

% Pengaturan Spasi (Spasi 1.5)
\setstretch{1.5}

% Pengaturan Paragraf (Indentasi alinea baru)
\setlength{\parindent}{1.25cm} % Sesuai "ketukan ke-7" (kira-kira 0.5 inci)
\setlength{\parskip}{0pt}

% Penomoran Bab (BAB I, BAB II, dst.)
\renewcommand{\thechapter}{\Roman{chapter}}

% Penomoran Section (A, B, C...) reset setiap Bab
\renewcommand{\thesection}{\Alph{section}}
\makeatletter
\@addtoreset{section}{chapter}
\makeatother

% Penomoran Subsection (1, 2, 3...) reset setiap Section
\renewcommand{\thesubsection}{\arabic{subsection}}
\makeatletter
\@addtoreset{subsection}{section}
\makeatother

% Penomoran Tabel, Gambar, Persamaan (1.1, 1.2, dst. - menggunakan angka Arab untuk bab)
\renewcommand{\thetable}{\arabic{chapter}.\arabic{table}}
\renewcommand{\thefigure}{\arabic{chapter}.\arabic{figure}}
\renewcommand{\theequation}{\arabic{chapter}.\arabic{equation}}

% Format Judul Bab (Level 0)
\titleformat{\chapter}[display]
  {\normalfont\bfseries\centering}
  {\MakeUppercase{\chaptername}\ \thechapter}
  {0pt}
  {\MakeUppercase}
\titlespacing*{\chapter}{0pt}{-20pt}{10pt}

% Format Judul Section (Level 1)
\titleformat{\section}
  {\normalfont\bfseries}
  {\thesection.\quad}
  {0pt}
  {}
\titlespacing*{\section}{0pt}{0pt}{0pt}

% Format Judul Subsection (Level 2)
\titleformat{\subsection}
  {\normalfont\bfseries}
  {\thesubsection.\quad}
  {0pt}
  {}

% Pengaturan Caption (Tabel di atas, Gambar di bawah)
\captionsetup{
  font=small,
  labelfont=bf,
  justification=centering,
  singlelinecheck=false,
  skip=10pt
}
\captionsetup[table]{position=above}
\captionsetup[figure]{position=below}

% --- Pengaturan Daftar Isi, Tabel, Gambar (TOC/LOT/LOF) ---
\setcounter{tocdepth}{1}
\setcounter{secnumdepth}{2}

% Menambahkan titik-titik (dotted leaders) untuk Bab
\renewcommand{\cftchapleader}{\cftdotfill{\cftdotsep}}

% Menambahkan kata "BAB" di depan nomor bab pada Daftar Isi
\renewcommand{\cftchappresnum}{BAB }
\setlength{\cftchapnumwidth}{1.8cm}
\renewcommand{\cftchapaftersnumb}{\quad}

% Membuat judul dertar isi, tabel, gambar berada di tengah dan All Caps
\renewcommand{\contentsname}{\centering DAFTAR ISI}
\renewcommand{\listtablename}{\centering DAFTAR TABEL}
\renewcommand{\listfigurename}{\centering DAFTAR GAMBAR}

% --- Pengaturan Posisi Nomor Halaman ---
\fancypagestyle{frontmatter}{
  \fancyhf{}
  \fancyfoot[C]{\thepage} % Tengah bawah untuk Romawi
  \renewcommand{\headrulewidth}{0pt}
}

% Redefine plain style specifically for frontmatter
\fancypagestyle{plain}{
  \fancyhf{}
  \fancyfoot[C]{\thepage}
  \renewcommand{\headrulewidth}{0pt}
}

\fancypagestyle{mainmatter}{
  \fancyhf{}
  \fancyfoot[R]{\thepage} % Kanan bawah untuk Arab
  \renewcommand{\headrulewidth}{0pt}
}

% Command to switch to mainmatter style including redefinition of plain
\newcommand{\startmainmatter}{
  \cleardoublepage
  \pagenumbering{arabic}
  \pagestyle{mainmatter}
  % Redefine plain style for chapter pages in mainmatter
  \fancypagestyle{plain}{
    \fancyhf{}
    \fancyfoot[R]{\thepage}
    \renewcommand{\headrulewidth}{0pt}
  }
}

% Helper command untuk istilah asing
\newcommand{\foreign}[1]{\textit{#1}}

% Persamaan matematis (bab.urutan) menggunakan angka Arab untuk bab
\numberwithin{equation}{chapter}
\renewcommand{\theequation}{\arabic{chapter}.\arabic{equation}}
